
\documentclass{article}
\usepackage{colortbl}
\usepackage{makecell}
\usepackage{multirow}
\usepackage{supertabular}

\begin{document}

\newcounter{utterance}

\centering \large Interaction Transcript for game `codenames', experiment `association\_difficult', episode 4 with reflexion{-}Qwen3.5{-}4B.
\vspace{24pt}

{ \footnotesize  \setcounter{utterance}{1}
\setlength{\tabcolsep}{0pt}
\begin{supertabular}{c@{$\;$}|p{.15\linewidth}@{}p{.15\linewidth}p{.15\linewidth}p{.15\linewidth}p{.15\linewidth}p{.15\linewidth}}
    \# & $\;$A & \multicolumn{4}{c}{Game Master} & $\;\:$B\\
    \hline

    \theutterance \stepcounter{utterance}  
    & & & \multicolumn{2}{p{0.3\linewidth}}{
        \cellcolor[rgb]{0.9,0.9,0.9}{
            \makecell[{{p{\linewidth}}}]{
                \texttt{\tiny{[GM$|$GM]}}
                {'hidden': {'team': ['polyester', 'gymnastics', 'nitrogen', 'hat'], 'innocent': ['ostrich', 'volleyball', 'korea', 'pants', 'archbishop'], 'opponent': [], 'assassin': ['coal']}, 'revealed': {'team': {'team': [], 'innocent': [], 'opponent': [], 'assassin': []}, 'opponent': {'team': [], 'innocent': [], 'opponent': [], 'assassin': []}}}
            }
        }
    }
    & & \\ \\

    \theutterance \stepcounter{utterance}  
    & & \multicolumn{4}{p{0.6\linewidth}}{
        \cellcolor[rgb]{0.9,0.9,0.9}{
            \makecell[{{p{\linewidth}}}]{
                \texttt{\tiny{[P1$\langle$GM]}}
                \texttt{You have attempted this same task before and did not succeed. These are the plans you wrote after those attempts.} \\
\texttt{Trial \#1: New plan: I failed to guess the final word "polyester" because I chose the clue "CLOTH" which was too generic and likely led to the opponent guessing it as a distractor or the partner failing to distinguish it from other cloth{-}related items before the turn ended. I should have selected a more specific or unique association for "polyester" that minimizes overlap with distractors like "wool" or "fabric." For the remaining words (polyester, gymnastics, nitrogen, hat), I need to prioritize clues that are distinctively related to only one or two specific targets to prevent the opponent from guessing them as distractors. Specifically, for "polyester," I should have used a clue like "PLASTIC" or "FIBER" (if safe) or "CHEMICAL" to differentiate it from natural fibers, ensuring the partner knows exactly which word to guess without ambiguity. I will also ensure I always list the targets clearly and avoid clues that could be misinterpreted as distractors.} \\
\\ 
\\ 
\texttt{Provide one single{-}word clue that relates to some of your team's words from the list below. You can choose to how many and to which words your clue relates to.} \\
\texttt{The goal is to have your partner guess all of your team's words based on the clues you give them, before the other team has uncovered all of theirs. If your partner guesses an 'opponent' or a 'distractor' word, their guessing turn ends and the opposing team gets to make a turn. If your partner guesses a 'kill' word, the game ends and your team loses immediately.} \\
\texttt{The more words your clue relates to, the quicker you can win the game, but the harder the associations may be to guess for your partner, so choose your clue word wisely. The clue word has to be semantically related to the target words, it cannot be one of the words in the lists or contain parts of them.} \\
\\ 
\texttt{Always give your single{-}word clue and your comma{-}separated list of related target words in the following format and make your answers as short as possible, never include any other text than is required in this form:} \\
\\ 
\texttt{CLUE: <WORD>} \\
\texttt{TARGETS: <TARGETS>} \\
\\ 
\texttt{Your team words are: polyester, gymnastics, nitrogen, hat.} \\
\texttt{Your opponent's team words are: .} \\
\texttt{Distractor words are: ostrich, volleyball, korea, pants, archbishop.} \\
\texttt{The kill words are: coal.} \\
            }
        }
    }
    & \\ \\

    \theutterance \stepcounter{utterance}  
    & \multicolumn{4}{p{0.6\linewidth}}{
        \cellcolor[rgb]{0.9,0.9,0.9}{
            \makecell[{{p{\linewidth}}}]{
                \texttt{\tiny{[P1$\rangle$GM]}}
                \texttt{CLUE: PLASTIC} \\
\texttt{TARGETS: polyester} \\
            }
        }
    }
    & & \\ \\

    \theutterance \stepcounter{utterance}  
    & & & \multicolumn{2}{p{0.3\linewidth}}{
        \cellcolor[rgb]{0.9,0.9,0.9}{
            \makecell[{{p{\linewidth}}}]{
                \texttt{\tiny{[GM$|$GM]}}
                \texttt{plastic} \\
            }
        }
    }
    & & \\ \\

    \theutterance \stepcounter{utterance}  
    & & & \multicolumn{2}{p{0.3\linewidth}}{
        \cellcolor[rgb]{0.9,0.9,0.9}{
            \makecell[{{p{\linewidth}}}]{
                \texttt{\tiny{[GM$|$GM]}}
                ['polyester']
            }
        }
    }
    & & \\ \\

    \theutterance \stepcounter{utterance}  
    & & & \multicolumn{2}{p{0.3\linewidth}}{
        \cellcolor[rgb]{0.9,0.9,0.9}{
            \makecell[{{p{\linewidth}}}]{
                \texttt{\tiny{[GM$|$GM]}}
                {'word': 'polyester', 'assignment': 'team'}
            }
        }
    }
    & & \\ \\

    \theutterance \stepcounter{utterance}  
    & & & \multicolumn{2}{p{0.3\linewidth}}{
        \cellcolor[rgb]{0.9,0.9,0.9}{
            \makecell[{{p{\linewidth}}}]{
                \texttt{\tiny{[GM$|$GM]}}
                \texttt{opponent has won} \\
            }
        }
    }
    & & \\ \\

    \theutterance \stepcounter{utterance}  
    & \multicolumn{4}{p{0.6\linewidth}}{
        \cellcolor[rgb]{0.95,0.95,0.95}{
            \makecell[{{p{\linewidth}}}]{
                \texttt{\tiny{[P2$\rightarrow$P1]}}
                \texttt{Post{-}episode Reflexion (not sent to GM):} \\
\texttt{New plan: I failed because "PLASTIC" is too broad and likely caused the partner to guess "polyester" too quickly or incorrectly, or perhaps the clue was invalid for the specific game rules regarding word lists. I need to select a clue that is uniquely associated with "polyester" but not with the distractors or other team words. A better strategy is to use a chemical or industrial term like "POLYMER" or "SYNTHETIC" if allowed, or focus on a specific property like "FIBER" (if not a distractor). However, looking at the distractors (ostrich, volleyball, korea, pants, archbishop), "PLASTIC" might have been valid but insufficient. The real issue might be that I didn't list all related targets or the clue was ambiguous. I will choose a more specific clue like "CHEMICAL" or "ARTIFICIAL" and ensure I list "polyester" as the target. Wait, "CHEMICAL" might be too generic. Let's try "FABRIC" if it's not a distractor, but "pants" is. Actually, "POLY" is part of the word, which is forbidden. "PLASTIC" was used and failed. Maybe the partner guessed a distractor? No, the loss was "failed to achieve target state". Perhaps the partner guessed "polyester" but it was too fast? Or maybe "PLASTIC" is considered a distractor word in some contexts? No, it's not in the list. The key is to make the clue so specific that only "polyester" fits. "PLASTIC" fits many things. I should use "POLYMER" or "SYNTHETIC". Let's go with "SYNTHETIC" as it is distinct from natural fibers and not in the distractor list. I will also ensure I list "polyester" clearly.} \\
\\ 
\texttt{New plan: I chose "PLASTIC" which is too common and might have led to the partner guessing a distractor or failing to distinguish it from other plastic items. I should have used a more specific term like "SYNTHETIC" or "POLYMER" that uniquely identifies "polyester" among the distractors and other team words. I will ensure the clue is semantically tight and list "polyester" as the target. I will also verify that the clue word is not a distractor or a kill word. Next time, I will use "SYNTHETIC" for "polyester" and list it as the target.} \\
            }
        }
    }
    & & \\ \\

\end{supertabular}
}

\end{document}
